\documentclass[../main.tex]{subfiles}
\begin{document}
%%%%%%%%%%%%%%%%%%%%%%%%%%%%%%%%%%%%%%%%%%%%%%%%%%%%%%%%%%%%%%%%%%%%%%%%%%%%%%%%%%%%%%%%%%%%%%%%%%%
\chapter{Using Data Variables}

In OmniGraffle, a \textit{data variable} is a string that you can use to insert object and line
labels based on other information in your project. These variables look for like-named
elements in an OmniGraffle project's XML source, and displays its corresponding data related
to the string provided. For example, you can use the \texttt{<\%Layer\%>} data variable as a
text object on the canvas to display the name of that layer.

Data variables have long been part of OmniGraffle; however, there is a little-known variable,
\texttt{<\%UserData\%>} which gives OmniGraffle Pro users access to project-level details and
metadata entered into the Object Data inspector.

This chapter provides you with a basic overview of the data variables available in OmniGraffle
Standard, and shows OmniGraffle Pro users how to go beyond the basics.

\section{Data Variable Basics}

You can insert a \textit{variable} while editing the label of an object or line, or a text
object. For example:

\begin{enumerate}
  \item Select the Text tool.
  \item Click on the canvas to create a text object.
  \item Enter the following:

    \texttt{<\%Canvas\%>}
  \item Press the Escape (\textbf{Esc}) key to accept the variable as the text.
\end{enumerate}

When you press the Escape key, the variable displays the name of the current canvas. If you
double-click on the text object, OmniGraffle displays the \texttt{<\%Canvas\%>} variable.

\ogfigure{0.800}{res/og7_variables01.png}

If you change the name of the canvas in the Sidebar (or in the Canvas Data inspector, if you
have OmniGraffle Pro), the variable instantly updates to reflect the new canvas name.

\ogfigure{0.800}{res/og7_variables02.png}

\subsection{The Core Set of Data Variables}

OmniGraffle Standard and Pro both have access to a core set of data variables, which you can
create without typing by choosing \textbf{Edit} $\blacktriangleright$ \textbf{Insert
Variable}, and then selecting one of the variable options from the submenu. This places a text
object on the current layer, containing the data variable of your choosing.

The core set of data variables include are listed in Table 1.

\begingroup
\small
\setlength{\LTleft}{0.075\textwidth}
\setlength{\LTright}{0.075\textwidth}
\renewcommand{\arraystretch}{1.25}
\begin{xltabular}{0.85\textwidth}{|l|>{\raggedright\arraybackslash}X|>{\raggedright\arraybackslash}X|}
  \caption{List of variables available from the Insert Variable option in the Edit menu}\\
  \hline
  \rowcolor{ogheaderblue}
  \textcolor{white}{\textbf{Menu Option}} & \textcolor{white}{\textbf{Variable}} &
  \textcolor{white}{\textbf{Description}} \\
  \hline
  \endhead
  \hline
  \endfoot
  \textbf{Canvas and Layer Variables} &  &  \\
  \hline
  Canvas Name & \ogcellcode{<\%Canvas\%>} & Displays the canvas name \\
  \hline
  Canvas Number & \ogcellcode{<\%CanvasNumber\%>} & Displays a number related to the canvas's
  position in the Sidebar \\
  \hline
  Total Canvases & \ogcellcode{<\%TotalCanvases\%>} & Displays the total number of canvases in
  the project \\
  \hline
  Layer Name & \ogcellcode{<\%Layer\%>} & Displays the name of the layer \\
  \hline
  \textbf{Document-specific Variables} &  &  \\
  \hline
  Document Name & \ogcellcode{<\%Document\%>} & Displays the filename of the OmniGraffle
  project \\
  \hline
  Creator & \ogcellcode{<\%Creator\%>} & Displays the user name of the person who created the
  project \\
  \hline
  Modifier & \ogcellcode{<\%Modifier\%>} & Displays the user name of the person who last
  modified the project \\
  \hline
  Page Number & \ogcellcode{<\%\#\%>} & Displays the current page number, based on the
  object's position on the canvas as it relates to how the page is printed \\
  \hline
  Total Pages & \ogcellcode{<\%TotalPages\%>} & Displays the total number of pages to be
  printed from the project \\
  \hline
  \textbf{Date Variables} &  &  \\
  \hline
  Current Date & \ogcellcode{<\%Date\%>} & Displays today's date \\
  \hline
  Document Creation Date & \ogcellcode{<\%CreationDate\%>} & Displays the date the project was
  created \\
  \hline
  Document Modification Date & \ogcellcode{<\%ModificationDate\%>} & Displays the date that
  the project was last modified \\
  \hline
  \textbf{Object-specific Variables} &  &  \\
  \hline
  X Position & \ogcellcode{<\%X\%>} & Displays the object's horizontal distance from the
  origin \\
  \hline
  Y Position & \ogcellcode{<\%Y\%>} & Displays the objects vertical distance from the origin
  \\
  \hline
  Width & \ogcellcode{<\%Width\%>} & Displays the object's width \\
  \hline
  Height & \ogcellcode{<\%Height\%>} & Displays the object's height \\
  \hline
  Line Length & \ogcellcode{<\%Length\%>} & Displays a line object's length \\
  \hline
  Graphic Name & \ogcellcode{<\%GraphicName\%>} & Displays the name of the object \\
  \hline
\end{xltabular}
\endgroup

\subsection{Customizing the Date Strings Used by the Date Variables}

As noted in Table 1, OmniGraffle has three date-specific variables that make it easy for you to
track the status of work on your projects.

\begin{itemize}
  \item \textbf{Current Date} (\texttt{<\%Date\%>}) --- inserts the current date. This
    particularly handy when you want to ``stamp'' a document you're exporting or printing with
    today's date.
  \item \textbf{Document Creation Date} (\texttt{<\%CreationDate\%>}) --- inserts the date that
    the document was originally created, according to the document XML.
  \item \textbf{Document Modification Date} (\texttt{<\%ModificationDate\%>}) --- inserts the
    date that the document was last saved, according to the document XML.
\end{itemize}

\begin{ogtip}
  The date variables are particularly useful when combined with the \texttt{<\%Creator\%>} and
  \texttt{<\%Modifier\%>} variables, when you need to see who last updated a project file.
\end{ogtip}

If you are using one of the date variables, you can customize the format of the date using the
variables listed in Table 2.

\begingroup
\small
\setlength{\LTleft}{0.075\textwidth}
\setlength{\LTright}{0.075\textwidth}
\renewcommand{\arraystretch}{1.25}
\begin{xltabular}{0.85\textwidth}{|l|>{\raggedright\arraybackslash}X|}
  \caption{Date and time variables}\\
  \hline
  \rowcolor{ogheaderblue}
  \textcolor{white}{\textbf{Variable}} & \textcolor{white}{\textbf{Description}} \\
  \hline
  \endhead
  \hline
  \endfoot
  \ogcellcode{y} & Two-digit year, based on the decade; for example, 16 for ``2016'' \\
  \hline
  \ogcellcode{Y} & Four-digit year; for example, 2016 \\
  \hline
  \ogcellcode{m} & Month, represented as two digits ranging from 01-12; for example, 10 for
  October \\
  \hline
  \ogcellcode{d} & Day, represented as two digits ranging from 01-31 (as appropriate); for
  example, 05 is the fifth day of the month \\
  \hline
  \ogcellcode{H} & Hours, represented as two-digits ranging from 00-23; for example, 11 is 11
  a.m., whereas 23 is 11 p.m. \\
  \hline
  \ogcellcode{M} & Minutes, represented as two digits ranging from 00-59 \\
  \hline
  \ogcellcode{S} & Seconds, represented as two digits ranging from 00-59 \\
  \hline
  \ogcellcode{z} & Zulu time difference; for example, an output of -0700 represents that the
  current time is seven hours behind Zulu time \\
  \hline
\end{xltabular}
\endgroup

For example, to get a date such as \texttt{2016-10-05 08:37:30 -0700Z}, you could use the
following variable string:

\texttt{<\%date \%Y-\%m-\%d \%H:\%M:\%S \%zZ \%>}

\begin{ognote}
  The logic we have used in creating the date variables is based on
  \href{http://developer.apple.com/library/ios/documentation/cocoa/Conceptual/DataFormatting/Articles/dfDateFormatting10_4.html}{Cocoa's
  date formatting tokens}.
\end{ognote}

\subsection{Using the Line Length Variable}

Apply the \textbf{Line Length} variable (\texttt{<\%Length\%>}) is especially useful when
assigned as the label on a line. The Line Length variable relies on the values set in the
\textbf{Units} inspector when displaying the line length.

For example, let's say that you are creating a floorplan of your house or apartment. In the
\textbf{Units} inspector, you can set the \textbf{Ruler Units} to \textbf{feet (ft)}, and the
\textbf{Scale} to something like \textbf{1 in = 1 ft} to get accurate (to-scale) floorplans,
using lines and the Line Length variable to display the dimensions.

\ogfigure{0.769}{res/og7_variables_linelength.png}

\section{Data Variables Bound to Inspectors}

Two of OmniGraffle's inspectors---Document Data (\textbf{Standard}), and Object Data
(\textbf{PRO})---offer a way for you to provide additional information about the objects in
your project, and the project file itself, with ease. And, you can tap into the information
you've provided in these inspectors with some special, inspector-bound variables.

\ogfigure{0.800}{res/og7_variables_docnoteinsp.png}

\begin{ognote}
  Similar in form to the Object Data inspector, the only information that you can access from the
  Canvas Data inspector is the \textbf{Name} field, via the \texttt{<\%Canvas\%>} variable.

  If you choose to add a Note or any key-value data to the Canvas Data inspector, that
  information is solely referential in nature. At present, OmniGraffle doesn't offer variables
  to access those fields.
\end{ognote}

\subsection{Variables for the Document Data Inspector}

Every field of the Document Data inspector has a like-named variable. This also includes the
metadata pop-up menu, which lets you add details about the project, including a list of
Authors, Organizations, Languages, Keywords, or Projects. And while the metadata you enter in
the rows beneath the pop-up menu can be singular, you can enter as many as you'd like.

The Document Data inspector's variables are not accessible via the \textbf{Edit}
$\blacktriangleright$ \textbf{Insert Variable} menu, and are shown and listed here:

\ogfigure{0.800}{res/og7_document_data_variables.png}

\begin{itemize}
  \item \texttt{<\%Subject\%>} --- Displays the Subject of the field.
  \item \texttt{<\%Copyright\%>} --- Books and software are protected by copyrights, and you can
    label your projects with a similar copyright statement.
  \item \texttt{<\%Version\%>} --- If needed, enter a version number for your project. This can
    be helpful for working on projects that are synced and shared among team members by having
    them iterate the version string whenever a change is made.
  \item \texttt{<\%Description\%>} --- Enter a brief description about the project.
  \item \texttt{<\%Comments\%>} --- Use the Comments field to enter a little more information
    about the project; more than you would enter in the Description field.
  \item \texttt{<\%Authors\%>} --- Here, ``Authors'' is used loosely to define the people
    associated with the project. The people you list in the rows underneath can be designers,
    engineers, architects, and so on.
  \item \texttt{<\%Organizations\%>} --- This could be the name of your company or the name of
    the client this project is related to.
  \item \texttt{<\%Languages\%>} --- Here you can specify the languages within the project. For
    example, if you are prototyping an app that will be localized in other languages, you can
    list those here.
  \item \texttt{<\%Keywords\%>} --- You can add metadata to your projects when Keywords is
    selected in the pop-up menu. There's no limit to the number of keywords you can assign to
    your project. When you use \texttt{<\%Keywords\%>}, each keyword you assign to a project is
    displayed on its own line.
  \item \texttt{<\%Projects\%>} --- Used for assigning a name to the project you are working on.
    You can assign more than one project name in the list; each one is displayed on its own line
    when the variable is called.
\end{itemize}

\subsection{Using the UserData Variable to Access Key-Value Pairs
(\texorpdfstring{\textbf{PRO}}{PRO})}

\begin{ogpro}
  If you have OmniGraffle Pro, you can enter details about specific objects in your projects using
  the Object Data inspector, and then---with the use of variables---apply that information to
  the objects on the canvas.

  The Object Data inspector looks fairly innocuous, there's the \textbf{Name} field at the top
  which is bound to the \texttt{<\%GraphicName\%>} variable. There's a \textbf{Note} field,
  where you can enter information about the object. The information you enter here is
  displayed whenever you hover over that object with the mouse. (There is no variable for
  accessing the Note field.)

  The \texttt{<\%GraphicName\%>} variable displays the name of the object, and is not dependent
  on having OmniGraffle Pro. If you have OmniGraffle Standard, you can assign names to objects
  in the Sidebar instead of in the Object Data inspector.

  To learn more about naming objects in the Sidebar, see Changing Canvas, Layer, and Object
  Names in the Sidebar, earlier in this guide.

  \subsubsection{Accessing Key-Value Pairs with the UserData Variable}

  At the bottom of the Object Data inspector is an area where you can assign \textit{key-value
  pairs} to individual objects in your project. A key-value pair is similar to the listings
  you find in a dictionary, where the word you look up is the \textit{key} and its definition
  is the \textit{value}. You can use the \texttt{<\%UserData\%>} variable to call upon the
  \textit{key} and use its \textit{value} as the label of an object or line.

  \ogfigure{0.800}{res/og7_variables_keyvalue01.png}

  Using the previous image as an example, \textbf{bucket02} is the \textit{key}, and
  \textbf{Article Bucket} is the \textit{value}. To display the value (Article Bucket) as the
  label for an object, you would use the following variable:

  \texttt{<\%UserData bucket02\%>}

  \ogfigure{0.800}{res/og7_variables_userdata_keyvalue.png}

  \subsubsection{Using Multiple Key-Value Pairs on the Same Object}

  You don't need to limit the number of key-value pairs you assign to an object to just one
  pair. For example, if you are working on a website design, you might want to devote a canvas
  in your project to hold detailed color swatches you plan to use on the site. You can then
  use the keys to identify color spaces, and the values to display the appropriate color
  value.

  \ogfigure{0.800}{res/og7_variables_keyvalue_colorswatches.png}

  \begin{description}[style=nextline]
    \item[{The variables used as the object labels are as follows:}] \texttt{<\%GraphicName\%>}
    \item[{}] \textbf{HEX:} \texttt{<\%UserData HEX\%>}
    \item[{}] \textbf{RGB:} \texttt{<\%UserData R\%>}\textbf{,}\texttt{<\%UserData
      G\%>}\textbf{,}\texttt{<\%UserData B\%>}
  \end{description}

  In this example, create a single object on the canvas, add the variables (shown above) as the
  text label for the object, and then use the Object Data inspector to assign the hexadecimal
  and RGB color values as key-value pairs. With those details in place, duplicate the swatch
  as many times as you need, and align them using the Smart Guides or with the Alignment
  inspector. Finally, use the Fill inspector to change the fill color for each object, and
  then apply the color values in the Object Data inspector.

  A similar technique could also be used to create a highly detailed network diagram. For
  example, you could note the IP addresses for servers, switches, and firewalls, the make and
  model of each piece of gear, the location of the equipment, and more.

  Look for ways to leverage key-value pairs and call upon them using the \texttt{<\%UserData
  [key]\%>} variable.
\end{ogpro}

\section{Chaining and Nesting Data Variables}

Another way you can use data variables in your projects is to combine them. For example, if you
want an object to display its width and height, you can set its label as \texttt{<\%Width\%> x
<\%Height\%>}, thus \textit{chaining} those two variables together.

Similarly, some variables allow you to \textit{nest} them together within a single set of angle
brackets. For example, you can combine the Page Number and Total Pages variables together as
follows to display the range of pages within your project:

\texttt{<\%\# of \%TotalPages \%>}

Not all variables, however, can be nested; for example, \texttt{<\%CanvasNumber of
\%TotalCanvases \%>} does not work. Instead, you'll have to chain those together as follows:

\texttt{<\%CanvasNumber\%> of <\%TotalCanvases\%>}

\section{Styling Variables}

The text produced by a variable can be styled similar to any other block of text by using the
\textbf{Font} and \textbf{Text Position} inspectors. You can change the font, type size, and
color in the Font inspector, and you can change the alignment, overflow state, and margins in
the Text Position inspector. The style changes you make in the inspectors are reflected in the
variable itself if you double- or triple-click on the text block containing the variable.

\ogfigure{0.800}{res/og7_variables_fontstyle.png}

\begin{ogpro}
  If you have OmniGraffle Pro, you can also convert variable-based text blocks to shapes, by
  selecting the text object and choosing \textbf{Edit} $\blacktriangleright$ \textbf{Objects}
  $\blacktriangleright$ \textbf{Convert Text to Shapes}.

  After converting the text block to shapes, the variable goes away, and OmniGraffle renders the
  text value provided by the variable as shapes.
\end{ogpro}

\section{Using Variables on Shared Layers (\texorpdfstring{\textbf{PRO}}{PRO})}

\begin{ogpro}
  Some variables, such as \texttt{<\%Canvas\%>}, \texttt{<\%CanvasNumber\%>},
  \texttt{<\%TotalCanvases\%>}, \texttt{<\%Layer\%>}, \texttt{<\%\#\%>} (page number), and
  \texttt{<\%TotalPages\%>} are useful for placing on a shared layer and then sharing that layer
  with every canvas in your project.

  For example, if your project has more than one canvas, you can use the following variable on a
  shared layer to display each canvas's number:

  \texttt{<\%CanvasNumber\%>}

  If you want to display the number of the current canvas in relation to the total number of
  canvases in your project (resulting in something like \textbf{Canvas 1 of 3}), use the
  following variable string:

  \texttt{<\%CanvasNumber\%> of <\%TotalCanvases\%>}

  If you wanted to add the canvas name into the mix, you could use the following:

  \texttt{<\%Canvas\%> (<\%CanvasNumber\%> of <\%TotalCanvases\%>)}

  Which would display something along the lines of \textbf{Canvas 1 (1 of 3)}, or whatever name
  you have assigned to the canvas.
\end{ogpro}
%%%%%%%%%%%%%%%%%%%%%%%%%%%%%%%%%%%%%%%%%%%%%%%%%%%%%%%%%%%%%%%%%%%%%%%%%%%%%%%%%%%%%%%%%%%%%%%%%%%
\end{document}
